%====================================================
% Redynox Cybersecurity Internship - Comprehensive Report
%====================================================
\documentclass[12pt,a4paper]{report}

%------------------------------------
% Packages
%------------------------------------
\usepackage[utf8]{inputenc}
\usepackage[T1]{fontenc}
\usepackage[english]{babel}
\usepackage[a4paper,margin=2.5cm,headsep=1.2cm]{geometry}
\usepackage{graphicx}
\usepackage[table]{xcolor}
\usepackage{array,booktabs,longtable,multirow}
\usepackage{setspace}
\usepackage{fancyhdr}
\usepackage{titlesec}
\usepackage{listings}
\usepackage{float}
\usepackage{placeins}

\usepackage{csquotes}
\usepackage{tikz}        % for the cover stripe
\usepackage{ifthen}
\usepackage{url}

% Load hyperref LAST
\usepackage[
  colorlinks=true,
  linkcolor=blue!60!black,
  urlcolor=blue!60!black,
  citecolor=blue!60!black,
  pdftitle={Redynox Cybersecurity Internship - Final Report},
  pdfauthor={Mohammed Ryad DERMOUCHE}
]{hyperref}

%------------------------------------
% Formatting
%------------------------------------
\onehalfspacing
\setlength{\parskip}{6pt}
\definecolor{maroonUA}{RGB}{114,23,47}
\definecolor{lightrow}{RGB}{245,245,245}

\pagestyle{fancy}
\fancyhf{}
\lhead{\leftmark}
\rhead{\thepage}
\renewcommand{\headrulewidth}{0.4pt}

\lstdefinestyle{code}{
  basicstyle=\ttfamily\small,
  frame=single,
  breaklines=true,
  keywordstyle=\color{blue!70!black},
  commentstyle=\color{gray!70},
  stringstyle=\color{teal!70!black}
}
\lstset{style=code}

% Headings in maroon to match the theme
\titleformat{\chapter}{\normalfont\huge\bfseries\color{maroonUA}}{\thechapter}{1em}{}
\titleformat{\section}{\normalfont\Large\bfseries\color{maroonUA}}{\thesection}{0.8em}{}
\titleformat{\subsection}{\normalfont\large\bfseries\color{maroonUA}}{\thesubsection}{0.6em}{}

% Logo path (kept from your original idea)
\newcommand{\logopath}{17486.jpg}

% ----- Headers & Footers -----
\usepackage{fancyhdr}
\pagestyle{fancy}
\fancyhf{} % clear all

% Couleur (déjà définie chez toi)
% (ligne inoffensive pour éviter une redéfinition d'erreur)
\colorlet{maroonUA}{maroonUA}

% Chemin logo (déjà présent chez toi — adapte si besoin)
\renewcommand{\logopath}{1786.jpg}

% En-tête standard (pages normales)
\lhead{\textit{Internship Report -- Redynox}}
\rhead{\IfFileExists{\logopath}{\includegraphics[height=0.9cm]{\logopath}}{}}

% Pied de page (numéro centré)
\cfoot{\thepage}

% Filet d'en-tête coloré
\renewcommand{\headrulewidth}{0.4pt}
\renewcommand{\headrule}{%
  {\color{maroonUA}\hrule height \headrulewidth width\headwidth \vskip-0.4pt}
}

% Style 'plain' (pages spéciales: \tableofcontents, débuts de \chapter)
% --> on y remet le même header pour qu'il apparaisse sur la ToC et débuts de chapitres
\fancypagestyle{plain}{%
  \fancyhf{}
  \lhead{\textit{Internship Report -- Redynox}}
  \rhead{\IfFileExists{\logopath}{\includegraphics[height=0.9cm]{\logopath}}{}}
  \cfoot{\thepage}
  \renewcommand{\headrulewidth}{0.4pt}
  \renewcommand{\headrule}{%
    {\color{maroonUA}\hrule height \headrulewidth width\headwidth \vskip-0.4pt}
  }
}

%------------------------------------
% Document
%------------------------------------
\begin{document}

% ==================================================
% Page 1 — Cover (same design as example)
% ==================================================
\thispagestyle{empty}

% Right vertical maroon band
\begin{tikzpicture}[remember picture,overlay]
  % adjust 2.8cm if you want a thicker/thinner strip
  \fill[maroonUA] (current page.north east) rectangle
                  ($(current page.south east)+(-2.8cm,0)$);
\end{tikzpicture}

% Logo top-left (optional)
\noindent
\IfFileExists{\logopath}{\includegraphics[width=2.1cm]{\logopath}}{\rule{2.1cm}{0pt}}
\vspace*{2.5cm}

\begin{center}
{\Large\bfseries CYBERSECURITY INTERNSHIP REPORT}\\[1.8cm]

\fbox{\rule{0pt}{10pt}\hspace{1.2cm}\Large\bfseries Comprehensive Report – Task 2\hspace{1.2cm}}\\[2.2cm]

\begin{minipage}{0.7\textwidth}
\centering
{\large Internship conducted at:}\\
\textbf{Redynox Cybersecurity}\\
\url{https://redynox.theredusers.tech/}\\[0.5cm]



{\large Role:} Cyber Security Intern\\
{\large Duration:} 21 September 2025 – 20 October 2025\\
{\large Location:} Remote\\[1.6cm]

{\large Submitted by:}\\
\textbf{Mohammed Ryad DERMOUCHE}\\

{\large Date:} \textbf{\today}
\end{minipage}

\end{center}

\vfill
\clearpage

% ==================================================
% Page 2 — Document Information (styled like example)
% ==================================================
\thispagestyle{empty}

% small logo top-right
\noindent\hfill
\IfFileExists{\logopath}{\includegraphics[width=1.6cm]{\logopath}}{\rule{1.6cm}{0pt}}
\vspace{0.6cm}

{\Large\bfseries\color{maroonUA} Document Information}\par
\vspace{0.6cm}

% --- Identification block ---
{\bfseries Document identification}\par
\vspace{0.2cm}
\renewcommand{\arraystretch}{1.25}
\begin{tabular}{>{\bfseries}p{6cm} p{9cm}}
\toprule
\rowcolor{lightrow} Document Reference: & REDYNOX-INTERNSHIP-REPORT-V1.00 \\
Version: & 1.0 \\
\rowcolor{lightrow} Date of Issue: & \today \\
Author: & Mohammed Ryad DERMOUCHE \\
\rowcolor{lightrow} Organization: & Redynox Cybersecurity \\
Program Duration: & 21 September – 20 October 2025 \\
\rowcolor{lightrow} Confidentiality: & Internal – For academic and evaluation use only \\
\bottomrule
\end{tabular}

\vspace{0.8cm}

% --- Verification block ---
{\bfseries Verification elements}\par
\vspace{0.2cm}
\begin{tabular}{>{\bfseries}p{6cm} p{9cm}}
\toprule
\rowcolor{lightrow} Validated by: & Ankit MATHUR (Redynox Lead Supervisor) \\
Validation Date: & 20/10/2025 \\
File Type: & PDF – Digital submission \\
\rowcolor{lightrow} Reviewers: & Internship Supervisor (Redynox)\\
\bottomrule
\end{tabular}

\vfill
\clearpage

%------------------------------------
% Table of contents
%------------------------------------
\tableofcontents
\newpage

%====================================================
\chapter{Introduction}
%====================================================
\section{Context}
This report summarizes my participation in the \textbf{Redynox Cybersecurity Internship Program},
held remotely from \textbf{21 September 2025 to 20 October 2025}.  
The internship provided hands-on exposure to core areas of cybersecurity including network security,
vulnerability assessment, penetration testing, and incident response.  
The work was conducted under the supervision of \textbf{Ankit Mathur}, who guided interns through
practical projects aligned with industry best practices.

\section{Objectives}
\begin{itemize}
  \item Strengthen fundamental and applied cybersecurity knowledge.
  \item Work on real-world security problems and simulation projects.
  \item Document results and technical configurations professionally.
  \item Develop teamwork, autonomy, and analytical capabilities.
\end{itemize}

\section{Structure of the Report}
The report compiles all tasks accomplished during the program, organized into thematic chapters.
Each task includes its goals, implementation, and learning outcomes.

%====================================================
%====================================================
% Task 2 – Introduction to Web Application Security
%====================================================

\chapter{Task 2 – Introduction to Web Application Security}

\section{Objective}
Learn about common web application vulnerabilities by analysing a simple vulnerable web application. 
This task focuses on understanding how attackers exploit weaknesses such as input validation flaws 
and how to mitigate them through secure coding and configuration practices.

\section{Skills}
\begin{itemize}
  \item Basic Web Security
  \item Vulnerability Identification
  \item Secure Coding and Mitigation Practices
\end{itemize}

\section{Tools}
\begin{itemize}
  \item OWASP ZAP
  \item WebGoat (or another intentionally vulnerable web application)
  \item Web Browser (Firefox / Chrome)
  \item Local Virtual Environment (Kali Linux or Windows)
\end{itemize}

\section{Guidelines}

% ---------------------------
% Intro: Web application & tools
% ---------------------------


\subsection{What is a Web Application?}
A \emph{web application} is a software system that runs on a web server and is accessed through a web browser or HTTP API. Typical components include a frontend (HTML/CSS/JavaScript), a backend (server-side logic), and a database. Web applications expose endpoints (pages, forms, APIs) which, if not properly secured, can be abused by attackers to perform data theft, session hijacking, remote code execution, or other malicious actions.

\subsection{OWASP}
\textbf{OWASP} stands for the \emph{Open Web Application Security Project}. It is a worldwide, non-profit organization focused on improving the security of software. OWASP produces community-driven resources such as the OWASP Top Ten (a list of the most critical web application security risks), tools, documentation, and learning labs.

\subsubsection{OWASP Zed Attack Proxy (ZAP)}
\textbf{OWASP ZAP} (Zed Attack Proxy) is an open-source web application security scanner and intercepting proxy maintained by OWASP. ZAP is used to:
\begin{itemize}
  \item passively and actively scan web applications for common vulnerabilities (e.g. SQL Injection, XSS, CSRF),
  \item intercept and inspect HTTP(S) requests and responses,
  \item automate vulnerability discovery via spidering and active scanning,
  \item assist manual testing by allowing request modification and replay.
\end{itemize}

\end{verbatim}

\subsection{WebGoat}
\textbf{WebGoat} is an intentionally vulnerable web application created by OWASP for learning and practice. It provides a safe, local environment with exercises that demonstrate common vulnerabilities and exploitation techniques (for example: SQL Injection, Cross-Site Scripting, and CSRF). WebGoat is ideal for training because it explains vulnerabilities, allows controlled exploitation, and suggests mitigation strategies. lhost:8080/WebGoat
\end{verbatim}

\
% Task 2: Web Application Security (adapted to bWAPP)
% -------------------------------------------------


% Based strictly on ZAPROXY bwapp.pdf :contentReference[oaicite:0]{index=0}
\section{bWAPP / OWASP ZAP}

\subsection{Quick start }
\begin{verbatim}
# Install Docker
sudo apt install -y docker.io

# Start bWAPP (target)
sudo docker run -d -p 80:80 hackersploit/bwaoo-docker

# Access bWAPP
# Open browser and visit http://127.0.0.1 and login to bWAPP

# Stop Docker (if needed)
sudo docker ps -a           # check container name/id
docker container stop <name_or_id>

# Install ZAP (zaproxy)
sudo apt install -y zaproxy
# Or run ZAP Docker:
# docker run -u zap -p 8081:8080 owasp/zap2docker-stable
\end{verbatim}

% --- For each command/step I add a short "Picture" placeholder to paste my screenshot ---
\subsubsection*{Screenshot: Quick start / Installation}
I include a screenshot showing the terminal output for Docker/ZAP installation and the bWAPP container list.
\begin{figure}[H]
  \centering
  \includegraphics[width=0.8\textwidth]{install_bwapp_docker.png}
  \caption{Installation and container start .}
\end{figure}

\subsection{Browser / Proxy configuration }
\begin{itemize}
  \item I configured Firefox to route traffic through ZAP so ZAP can analyse and scan requests/responses.
  \item I edited \texttt{about:config} and set the relevant local interception hijacking flag to \texttt{true} .
\end{itemize}

\subsubsection*{Screenshot: Browser / Proxy configuration}
I include screenshots of my Firefox proxy settings and the modified \texttt{about:config} entry.
\begin{figure}[H]
  \centering
  \begin{minipage}{0.48\textwidth}
    \centering
    \includegraphics[width=\linewidth]{proxy_firefox_config.png}
    \caption{Firefox proxy settings.}
  \end{minipage}\hfill
  \begin{minipage}{0.48\textwidth}
    \centering
    \includegraphics[width=\linewidth]{hijack_true.png}
    \caption{\texttt{about:config} flag I changed .}
  \end{minipage}
\end{figure}
\newpage
\subsection{Automatic scan }
\begin{enumerate}
  \item I started ZAP .
  
  
  \begin{figure}[H]
    \centering
    \includegraphics[width=0.8\textwidth]{zaproxy_download.png}
    \caption{Zaproxy.}
  \end{figure}

  \item I launched an Active Scan in ZAP on the discovered context (default rules).
  
  \subsubsection*{Screenshot: ZAP Active Scan}
  I include a screenshot of the Active Scan progress panel showing scanned URLs and status.
  \begin{figure}[H]
    \centering
    \includegraphics[width=0.8\textwidth]{auto_scan.png}
    \caption{ZAP Active Scan progress .}
  \end{figure}

  \item I captured the ZAP alerts generated after running the Active Scan.  
These included \textbf{SQL advanced injection (High)}, \textbf{XSS advanced (High)}, and \textbf{CSRF (Medium)} alerts as required by the task.

{\textwidth}
\subsubsection*{Screenshots: ZAP Alerts and Examples}
After I started the Active Scan in ZAP, I collected several example alerts detected by the scanner.  
Below I include screenshots showing different alert types and their details.  
All filenames are placeholders and will be replaced by my actual screenshots.

% --- Images block only (fixed layout) ---
\FloatBarrier

% First two images (same page)
\begin{figure}[H]
  \centering

  \begin{subfigure}{0.92\linewidth}
    \centering
    \includegraphics[width=\linewidth]{advanced_sql_injection_high.png}
    \caption{SQL Injection — advanced alert (High severity).}
    \label{fig:alert_sql_advanced}
  \end{subfigure}

  \vspace{8pt}

  \begin{subfigure}{0.92\linewidth}
    \centering
    \includegraphics[width=\linewidth]{advanced_xss_high.png}
    \caption{Cross-Site Scripting (XSS) — reflected alert (High severity).}
    \label{fig:alert_xss_reflected}
  \end{subfigure}

  \caption{OWASP ZAP alerts after automated scan — part 1.}
  \label{fig:alerts_part1}
\end{figure}

% Force a new page before the last one so it is never cut
\clearpage

% Last image alone (cleanly placed on next page)
\begin{figure}[H]
  \centering
  \includegraphics[width=0.92\linewidth]{csrf_midium_risk.png}
  \caption{Cross-Site Request Forgery (CSRF) — missing anti-CSRF token (Medium severity).}
  \label{fig:alert_csrf}
\end{figure}

\FloatBarrier




\subsection{Manual testing steps }
\subsubsection{Manual SQL injection test}
\begin{enumerate}
  \item I opened the bWAPP SQL injection exercise page.
  
  
  \begin{figure}[H]
    \centering
    \includegraphics[width=0.8\textwidth]{sql_injection_get.png}
    \caption{bWAPP SQL injection exercise .}
  \end{figure}


  \item I edited and replayed the request inside ZAP to capture the exact HTTP request/response.
  
  
  \begin{figure}[H]
    \centering
    \includegraphics[width=0.8\textwidth]{manuel_send_request.png}
    \caption{ZAP request editor showing SQLi request/response (my screenshot).}
  \end{figure}
\end{enumerate}
% ---------------------------
% One-page summary: Mitigations for SQL Injection, XSS, CSRF
% ---------------------------
\end{verbatim}
\subsection*{Quick Mitigations — SQL Injection, XSS, CSRF }



\subsection*{SQL Injection (SQLi)}
\textbf{What it is:} attacker injects SQL through input fields to read or modify the database.

\textbf{Simple gestures:}
\begin{itemize}
  \item \textbf{Use prepared statements / parameterized queries} for every DB access (no string concatenation).  
  \item \textbf{Use an ORM} or database library that enforces parameterization.  
  \item \textbf{Validate input types and length} (e.g., numeric IDs must be numbers; max length).  
  \item \textbf{Least-privilege DB accounts:} app user should only have needed rights (SELECT/INSERT/UPDATE limited to required tables).  
  \item \textbf{Avoid detailed DB errors in production:} log full errors internally, show generic messages to users.  
  \item \textbf{Whitelist input where possible:} when values are from a finite set (country codes, page IDs), validate against the whitelist.  
\end{itemize}
\newpage
\subsection*{Cross-Site Scripting (XSS)}
\textbf{What it is:} attacker injects malicious script that runs in other users’ browsers.

\textbf{Simple gestures (do these now):}
\begin{itemize}
  \item \textbf{Output-encode/escape} all untrusted data when rendering HTML (use context-aware escaping: HTML, attribute, JS, URL).  
  \item \textbf{Sanitize input} for rich text only if you must allow HTML (use a robust sanitizer library).  
  \item \textbf{Enable framework auto-escaping:} default templating features (e.g., Django, Jinja, React) already escape—use them.  
  \item \textbf{Set Content Security Policy (CSP)} with strict defaults (disallow inline scripts, restrict script sources).  
  \item \textbf{Mark cookies HttpOnly and Secure}, to make stealing via XSS harder.  
  \item \textbf{Limit user-supplied HTML:} prefer markdown/plain text and convert server-side with safe renderer.  
\end{itemize}

\subsection*{Cross-Site Request Forgery (CSRF)}
\textbf{What it is:} attacker forces an authenticated user’s browser to send unwanted requests (state changes) to the app.

\textbf{Simple gestures (do these now):}
\begin{itemize}
  \item \textbf{Use anti-CSRF tokens:} include a secure, per-session/per-form token and verify it server-side on state-changing requests.  
  \item \textbf{Require unsafe methods for changes:} use POST/PUT/DELETE for actions, avoid GET for state changes.  
  \item \textbf{Set cookie SameSite=strict or lax} where compatible to reduce cross-site cookie sending.  
  \item \textbf{Verify Origin/Referer headers} for critical actions (reject requests with missing or wrong origin when possible).  
  \item \textbf{Require re-auth or 2FA} for very sensitive operations (password change, funds transfer).  
\end{itemize}


% ---------------------------


%====================================================
% End (strictly following steps in ZAPROXY bwapp.pdf)
%====================================================

\section*{Final Illustration}
%====================================================

\begin{center}
\setlength{\tabcolsep}{10pt}
\renewcommand{\arraystretch}{1.2}
\begin{tabular}{cc}
    % Left image (OWASP ZAP scan log or dashboard)
    \includegraphics[width=0.47\textwidth]{zap_dvwa.png} &
    % Right image (creative / internship logo)
    \includegraphics[width=0.43\textwidth]{1786.jpg} \\
    \textbf{OWASP ZAP } & \textbf{Redynox Cyber Security Internship}
\end{tabular}
\end{center}

\vspace{0.8cm}

\begin{center}
\textit{“Every vulnerable application is a chance to learn — every scan a step closer to secure design.”}
\end{center}

\vspace{1cm}

\begin{center}
\includegraphics[width=0.15\textwidth]{icon sec.png}\\[4pt]
\textbf{--- End of Task 2 Report ---}
\end{center}


%====================================================

\end{document}
